\section{The actual Wach coefficient embedding}
\label{sec:wach-inclusion-saturation-v054}

\status{This chapter proves all four clauses of the original
Wach-inclusion theorem: injectivity, reflection of units, saturation
for multiplication by the residue characteristic, and the exact
intersection with the integral Witt tensor ring. The supplied
coefficient series ring, period-ring maps and Frobenius are kept
throughout. This supplies one of the three period-ring inputs
retained in WLD26. The actual Wach prism and the invariant
Frobenius-matrix descent remain separate inputs.}

The coefficient embedding is not an arbitrary map between two fields.
Its source is a completed Laurent-series ring with several coefficient
components. Its target is an integral coefficient extension of a Witt
ring. Neither target is assumed to be a domain after tensoring.
The key observation is that a nonzero component idempotent cannot be
divisible by an element of the Jacobson radical. Together with the
first coefficient of a power series not divisible by a uniformizer,
this reflects divisibility from the Witt tensor ring back to the
given series ring.

We first establish the actual ring properties used in this argument.
The proofs use standard elementary foundations of finite extensions
of $\mathbb Q_p$, tensor products, and the coordinate formulas for
$p$-typical Witt vectors. No Wach module, period comparison functor,
or invariant Frobenius-solution descent is an input.

\subsection{The actual coefficient and Witt tensor rings}

% BEGIN WIS1
\paragraph{WIS1: complete integral coefficients and one uniformizer.}
\label{claim:WIS1}
Let $L/\mathbb Q_p$ be a finite extension and let $O$ be the integral
closure of $\mathbb Z_p$ in $L$. Then $O$ is a complete discrete
valuation ring, finite free of positive rank over $\mathbb Z_p$,
and its fraction field is $L$. There are a uniformizer $\varpi$,
an integer $e\geq1$, and $u\in O^\times$ such that
\[
                         p=u\varpi^e.                     \tag{WIS.1}
\]
Both the $p$-adic and $\varpi$-adic topologies on $O$ are complete
and separated.

\begin{proof}
Use the unique extension of the $p$-adic absolute value to $L$.
An element is integral over $\mathbb Z_p$ exactly when its absolute
value is at most one. Indeed, a monic integral equation precludes
absolute value greater than one by the ultrametric inequality.
Conversely, all conjugates have the same absolute value, so the
coefficients of the minimal polynomial have absolute value at most
one and belong to $\mathbb Z_p$. An integral element is a unit of
$O$ exactly when its absolute value is one: its inverse is integral
in that case, and the two integral factors of a product equal to
one must each have absolute value one. Thus the nonunits form
the ideal of elements with absolute value less than one, and $O$
is local.

For completeness we recall why the integral closure is finite.
Choose a $\mathbb Q_p$-basis $b_1,\ldots,b_d$ of $L$ consisting of
integral elements, after scaling an arbitrary basis. The trace
pairing is nondegenerate, since this characteristic-zero extension
is separable. For $x\in O$, each $\operatorname{Tr}_{L/\mathbb Q_p}
(xb_i)$ belongs to $\mathbb Z_p$: the trace is integral and lies
in $\mathbb Q_p$. Hence $O$ lies in the trace-dual lattice of
$\sum_i\mathbb Z_p b_i$. It is a submodule of a finite free
$\mathbb Z_p$-module and is therefore finite. It is torsion-free,
so the elementary structure theorem over the principal ideal
domain $\mathbb Z_p$ makes it finite free; its rank is $d>0$.
Scaling any element of $L$ by a sufficiently large power of $p$
puts it in $O$, proving the fraction-field assertion.

The finite integral closure is a one-dimensional integrally closed
Noetherian domain. It is local and not a field, because $p$ is
nonzero and has absolute value less than one. The discrete
valuation criterion for a one-dimensional normal Noetherian local
domain gives a uniformizer $\varpi$. Factoring the nonzero nonunit
$p$ gives (WIS.1). A finite free $\mathbb Z_p$-module is $p$-adically
complete and separated by taking limits in each basis coordinate.
The identity $pO=\varpi^eO$ makes the two stated filtrations cofinal,
which proves the remaining completeness assertions.
\end{proof}
% END WIS1

% BEGIN WIS2
\paragraph{WIS2: the Witt-coordinate properties used below.}
\label{claim:WIS2}
Let $S$ be a perfect ring of characteristic $p$. The ring $W(S)$
is $p$-adically complete and separated, multiplication by $p$ is
injective, and the constant-coordinate map has kernel $pW(S)$.
Moreover,
\[
       a\in W(S)^\times\quad\Longleftrightarrow\quad
                         a_0\in S^\times.                 \tag{WIS.2}
\]
For an inclusion of perfect characteristic-$p$ rings $S\hookrightarrow F$,
the induced map $W(S)\to W(F)$ is injective and its image is
$p$-saturated.

\begin{proof}
In Witt coordinates the multiplication formula is
\[
          (pa)_0=0,\qquad (pa)_{r+1}=a_r^p.                 \tag{WIS.3}
\]
Frobenius on $S$ is bijective. Thus a vector is divisible by $p$
exactly when its zeroth coordinate vanishes, and division by $p$
is unique. Iteration says that divisibility by $p^N$ is equivalent
to vanishing of the first $N$ coordinates. In particular,
multiplication by $p$ is injective and the $p$-adic filtration is
separated.

A compatible system modulo $p^N$ determines each of its coordinates
and hence one Witt vector. Addition and multiplication on every
finite truncation are given by the Witt polynomials; the vector
so obtained has the prescribed truncations. This proves
completeness. In particular $pW(S)$ lies in the Jacobson radical:
for any $b$, the series $\sum_{r\geq0}(pb)^r$ is an inverse of
$1-pb$.

Units map to units under the zeroth-coordinate homomorphism.
Conversely, if $a_0$ is a unit, multiply $a$ by the inverse of
the Teichmuller vector $[a_0]$. The result is $1+pb$ for some
$b$, and the preceding geometric series makes it a unit.
This proves (WIS.2).

An inclusion of coefficient rings gives an inclusion coordinate
by coordinate. If the image of $a\in W(S)$ is $pw$ in $W(F)$,
then (WIS.3) gives $w_r^p=a_{r+1}$ in $F$ for every $r$.
Let $b_r$ be the unique $p$th root of $a_{r+1}$ in $S$.
Injectivity of Frobenius in $F$ gives $w_r=b_r$, so
$w$ is the image of the Witt vector $(b_r)_{r\geq0}$ in $W(S)$.
This is the asserted saturation, including the actual divided
element rather than only its existence in a quotient.
\end{proof}
% END WIS2

% BEGIN WIS3
\paragraph{WIS3: the cyclotomic element is in the maximal ideal of $A_{\inf}$.}
\label{claim:WIS3}
Use the original tilt ring
\[
 S=\mathcal O_C^\flat
   =\varprojlim_{x\mapsto x^p}(\mathcal O_C/p),\qquad
 A_{\inf}=W(S).
\]
Then $S$ and $A_{\inf}$ are local. For the given compatible
cyclotomic system $\epsilon$ and its $p^s$th root $\epsilon_s$,
the element
\[
                         \mu_s=[\epsilon_s]-1             \tag{WIS.4}
\]
belongs to the maximal ideal of $A_{\inf}$ for every $s\geq0$.

\begin{proof}
The ring $\mathcal O_C$ is a valuation ring and $p$ is a nonunit,
so $\mathcal O_C/p$ is a nonzero local ring. For a compatible
sequence $x=(x_r)$ in its perfection, one has
$x_r^{p^r}=x_0$. If $x_0$ is a unit, every $x_r$ is a unit,
and the inverses form a compatible sequence. Thus $x$ is a unit
exactly when $x_0$ is a unit. Its nonunits form the inverse
image of the maximal ideal under the zeroth projection, proving
locality of $S$. WIS2 proves locality of $W(S)$ in the same way.

The original cyclotomic data give
$\operatorname{pr}_0(\epsilon)=\zeta_0\bmod p=1$.
Hence $\epsilon-1$ is not a unit.
In characteristic $p$,
\[
                     (\epsilon_s-1)^{p^s}=\epsilon-1.
\]
A power of a unit would be a unit, so $\epsilon_s-1$ is not
a unit. It is the zeroth Witt coordinate of $\mu_s$.
By (WIS.2), $\mu_s$ is a nonunit of the local ring $A_{\inf}$,
as required. This also covers $s=0$; no nonzero-coordinate
assumption is inserted.
\end{proof}
% END WIS3

% BEGIN WIS4
\paragraph{WIS4: finite coefficient coordinates commute with the actual maps.}
\label{claim:WIS4}
Fix a $\mathbb Z_p$-basis $b_1,\ldots,b_d$ of $O$. For every
$\mathbb Z_p$-algebra $D$, each element of $O\otimes_{\mathbb Z_p}D$
has a unique expression
\[
                          \sum_{i=1}^d b_i\otimes d_i.      \tag{WIS.5}
\]
Coefficient maps in these coordinates commute with maps $D\to D'$
and with multiplication by $1\otimes z$. In particular,
\[
 x\in(1\otimes z)^N(O\otimes D)
       \quad\Longleftrightarrow\quad d_i\in z^ND
                      \text{ for every }i.               \tag{WIS.6}
\]

\begin{proof}
Tensor the basis isomorphism $O\simeq\mathbb Z_p^d$ with $D$.
Its inverse and coordinate maps send $o\otimes a$ to the
coordinates of $o$ multiplied by $a$. These formulas prove
the two compatibility assertions first on pure tensors and
then by additivity. If $x=(1\otimes z)^Ny$, its coordinates
are $z^N$ times those of $y$. Conversely, choose a quotient
for each coordinate and use their finite sum (WIS.5).
This proves (WIS.6) without assuming that $D$ is flat.
\end{proof}
% END WIS4

% BEGIN WIS5
\paragraph{WIS5: the actual integral inclusion is injective and saturated.}
\label{claim:WIS5}
Let $F=\operatorname{Frac}(S)$ be the original tilt field and set
\[
 B=O\otimes_{\mathbb Z_p}A_{\inf},\qquad
 W=O\otimes_{\mathbb Z_p}W(F),\qquad c:B\longrightarrow W.
                                                               \tag{WIS.7}
\]
Both rings are $p$-adically complete and separated, multiplication
by $p$ is injective on them, $W\ne0$, and $c$ is injective.
Its image is $p$-saturated:
\[
                  c(b)=pw\ \Longrightarrow\
                        w=c(b')\text{ for some }b'\in B.   \tag{WIS.8}
\]

\begin{proof}
The supplied tilt construction makes $S$ a perfect integral
domain. Its fraction field is perfect as well, by taking
$p$th roots in the numerator and denominator. Apply WIS2
to $S$ and $F$, and then WIS4 to the fixed integral basis
of $O$. Completeness, separatedness, absence of $p$-torsion,
and injectivity of $c$ follow in each coordinate.
The underlying $W(F)$-module of $W$ is $W(F)^d$ with
$d>0$, so it is nonzero. For (WIS.8), each coordinate
of $c(b)=pw$ lies in the $p$-saturated image of
$A_{\inf}\to W(F)$ by WIS2. The finitely many divided
coordinates assemble to $b'$ using (WIS.5).
\end{proof}
% END WIS5

% BEGIN WIS6
\paragraph{WIS6: a finite matrix proves the required Jacobson-radical property.}
\label{claim:WIS6}
If $D$ is a commutative ring, $z\in\operatorname{Jac}(D)$,
and $C$ is a finite free commutative $D$-algebra, then
the image of $z$ lies in $\operatorname{Jac}(C)$. In particular,
for the rings (WIS.7),
\[
             \pi_B=1\otimes\mu_s\in\operatorname{Jac}(B),
             \qquad p\in\operatorname{Jac}(W).             \tag{WIS.9}
\]

\begin{proof}
For $b\in C$, multiplication by $1-zb$ on a fixed $D$-basis
has matrix $1-zM_b$. Its determinant is $1+zq$ for some
$q\in D$, and is a unit because $z$ belongs to the Jacobson
radical. The adjugate matrix gives an inverse linear map.
Its value at $1$ is an element $v\in C$ satisfying
$(1-zb)v=1$. Thus $1-zb$ is a unit for every $b$, which
is exactly the Jacobson-radical criterion.

The $A_{\inf}$-algebra $B$ is finite free by WIS4, and
$\mu_s\in\operatorname{Jac}(A_{\inf})$ by WIS3, proving
the first assertion in (WIS.9). The second follows directly
from $p$-adic completeness of $W$: the geometric series
in $pw$ inverts $1-pw$ for every $w\in W$.

Only (WIS.9) is needed in the embedding proof. This argument
does not assert that $B$ is a domain or local, and does not
assert the stronger $\mu_s$-adic completeness theorem.
\end{proof}
% END WIS6

\subsection{Power-series components and the abstract embedding problem}

% BEGIN WIS7
\paragraph{WIS7: the first coefficient not divisible by a uniformizer.}
\label{claim:WIS7}
Let $O$ be a local ring in which every nonunit is divisible by
$\varpi$. For each $a\in O[[T]]$, either $a=\varpi b$,
or there are $b,u\in O[[T]]$ and $k\geq0$ with
\[
                    u\in O[[T]]^\times,\qquad
                    a=\varpi b+T^ku.                      \tag{WIS.10}
\]
For $A=\prod_{i\in J}O[[T]]$, put
$t=(\varpi)_i$, $\pi=(T)_i$, and let $e_i$ be the component
idempotents. For every $a\in A$ and $i\in J$, either
$e_i a=tb$, or
\[
                     e_i a=tb+\pi^kue_i,\qquad u\in A^\times.
                                                               \tag{WIS.11}
\]

\begin{proof}
If every coefficient is divisible by $\varpi$, divide
the coefficients individually. Otherwise let $k$ be the
smallest index of a coefficient not divisible by $\varpi$.
Divide the coefficients below $k$ by $\varpi$ to form a
polynomial $b$, and take
$u=\sum_{r\geq0}a_{k+r}T^r$. Its constant coefficient
is a unit. The usual coefficient recursion constructs an
inverse power series for $u$, proving (WIS.10).

Apply this construction to $a_i$. Place $b$ in its $i$th
component and zero elsewhere. For the second alternative,
take the same unit series $u$ in every component; its product
with $e_i$ has only the needed component. This proves
(WIS.11), with a unit of the entire product ring.
\end{proof}
% END WIS7

% BEGIN WIS8
\paragraph{WIS8: coefficientwise completeness of the integral series ring.}
\label{claim:WIS8}
For the complete ring $O$ of WIS1, the ring
$A=\prod_{i\in J}O[[T]]$ is $p$-adically complete and separated.
In particular, if $a_k\in A$ and
\[
           a_k-a_{k'}\in p^kA\qquad(k\leq k'),
                                                               \tag{WIS.12}
\]
there exists $a\in A$ such that $a_k-a\in p^kA$ for every $k$.

\begin{proof}
For every component $i$ and degree $r$, take the $p$-adic
limit in $O$ of the coefficient of $T^r$ in $a_k(i)$.
These limits define a power series in each component.
Reducing modulo $p^k$ shows that the difference from $a_k$
has every coefficient divisible by $p^k$. Choosing those
quotients coefficientwise proves divisibility in $A$ itself.
If an element lies in every $p^kA$, each coefficient lies
in $\bigcap_kp^kO=0$, giving separatedness. The argument
even works for an arbitrary index set $J$; later its
finiteness is used to sum the component idempotents.
\end{proof}
% END WIS8

% BEGIN WIS9
\paragraph{WIS9: Frobenius keeps every component visible.}
\label{claim:WIS9}
Let $A=\prod_{i\in\mathbb Z/n\mathbb Z}O[[T]]$ with $n>0$,
let $j:A\to E$ and $f:E\to W$ be unital ring maps with
$W\ne0$, and suppose operators $\phi_A,\phi_E,\phi_W$ satisfy
\[
 \phi_A(e_i)=e_{i+1},\qquad
 \phi_Ej=j\phi_A,\qquad f\phi_E=\phi_Wf,\qquad \phi_W(0)=0.
                                                               \tag{WIS.13}
\]
Then $f(j(e_i))\ne0$ for every $i$.

\begin{proof}
If one image were zero, iterating (WIS.13) would make
the images of $e_i,e_{i+1},e_{i+2},\ldots$ all zero.
These run through all components. But the idempotents sum
to $1$, whose image is $1\ne0$ in $W$.
For the original series Frobenius,
$\phi_A(a)_i=\phi_0(a_{i-1})$, the first identity in
(WIS.13) follows because $\phi_0$ preserves zero and one.
No injectivity or surjectivity of a Frobenius operator
is needed for this argument.
\end{proof}
% END WIS9

We now formulate the algebraic problem in a way that keeps
the four original maps distinct. Let
\[
 A\xrightarrow{j}E\xrightarrow{f}W,\qquad
 A\xrightarrow{g}B\xrightarrow{c}W,\qquad fj=cg              \tag{WIS.14}
\]
be a commutative square of unital maps. In the next eight claims
we use precisely the following data.

\begin{enumerate}
\item In $A$ there are $t,\pi,u$ with $u$ a unit and
      $p=ut^e$ for an integer $e\geq1$.
\item A finite family of idempotents $e_i$ sums to one,
      every $f(j(e_i))$ is nonzero, and the dichotomy
      (WIS.11) holds for every $a$ and $i$.
\item $A$ has the completeness property (WIS.12).
      The ring $E$ is $p$-adically complete and separated,
      and $j(\pi)$ is a unit.
\item For every $x\in E$ there are $a\in A$, $m\geq0$ such that
      $xj(\pi)^m-j(a)\in pE$. Also, for every $a\in A$ and $k$,
      $j(a)\in p^kE$ implies $a\in p^kA$.
\item In $W$, multiplication by $p$ is injective and
      $p\in\operatorname{Jac}(W)$. The image of $c$ is
      $p$-saturated, in the sense of (WIS.8), and
      $g(\pi)\in\operatorname{Jac}(B)$.
\end{enumerate}

Write $t_E=j(t)$, $\pi_E=j(\pi)$, and
$t_W=f(t_E)$, $\pi_W=f(\pi_E)$; a subscript on another
element of $A$ denotes its image by the same maps.
At this stage neither $f$ nor $g$ is assumed injective.
For a subring image, notation such as $c(B)$ means its
set-theoretic image.

% BEGIN WIS10
\paragraph{WIS10: the uniformizer is radical and regular where needed.}
\label{claim:WIS10}
Under (WIS.14) and the listed hypotheses,
\[
 t_E\in\operatorname{Jac}(E),\qquad
 t_W\in\operatorname{Jac}(W),\qquad
 t_Ww=0\Longrightarrow w=0.                                \tag{WIS.15}
\]
No nonzero idempotent $f(j(e_i))$ is divisible by $t_W$.
Every $x\in E$ has a representation
\[
                     x\pi_E^m=j(a)+t_Ez                   \tag{WIS.16}
\]
for some $a\in A$, $z\in E$, $m\geq0$.

\begin{proof}
Completeness puts $p$ in $\operatorname{Jac}(E)$ by the
geometric series argument. Since $p=u_Et_E^e$ and $u_E$
is a unit, $t_E^e$ belongs to every maximal ideal.
Each maximal ideal is prime, hence contains $t_E$.
The same proof in $W$ uses the supplied radical property
of $p$. If $t_Ww=0$, multiplying by $u_Wt_W^{e-1}$
gives $pw=0$, so $w=0$.

An idempotent $h$ in the Jacobson radical is zero:
$1-h$ is a unit and $h(1-h)=0$. Thus if $t_W$ divided
$f(j(e_i))$, that nonzero idempotent would be zero.
Finally, the supplied density condition writes
$x\pi_E^m-j(a)=pr=u_Et_E^er$.
Take $z=u_Et_E^{e-1}r$ to obtain (WIS.16).
\end{proof}
% END WIS10

% BEGIN WIS11
\paragraph{WIS11: reflection of one uniformizer factor.}
\label{claim:WIS11}
For every $x\in E$,
\[
                         t_W\mid f(x)
                    \quad\Longrightarrow\quad t_E\mid x. \tag{WIS.17}
\]

\begin{proof}
Choose (WIS.16). Its image in $W$ shows that $t_W$
divides $f(j(a))$. Fix a component $i$. If the second
alternative of (WIS.11) held, mapping it to $W$ would give
\[
       \pi_W^kv_W f(j(e_i))
          =f(j(e_i))f(j(a))-t_Wf(j(b)).
\]
The right side is divisible by $t_W$, while
$\pi_W^kv_W$ is a unit. This would make the nonzero
idempotent $f(j(e_i))$ divisible by $t_W$, contrary
to WIS10. Thus $e_i a$ is divisible by $t$ for every $i$.
Summing over the finite family gives $a=tb$. Equation
(WIS.16) becomes $x\pi_E^m=t_E(j(b)+z)$.
Cancel the unit $\pi_E^m$ to obtain (WIS.17).
\end{proof}
% END WIS11

% BEGIN WIS12
\paragraph{WIS12: reflection of every uniformizer power.}
\label{claim:WIS12}
For every $k\geq0$ and $x\in E$,
\[
                      t_W^k\mid f(x)
                 \quad\Longrightarrow\quad t_E^k\mid x.   \tag{WIS.18}
\]

\begin{proof}
The case $k=0$ is immediate. If
$f(x)=t_W^{k+1}y$, WIS11 gives $x=t_Ex'$.
Therefore
$t_W(f(x')-t_W^ky)=0$. By the regularity in WIS10,
$f(x')=t_W^ky$. Apply the induction hypothesis to $x'$
and multiply its resulting divisibility by $t_E$.
The cancellation takes place in $W$ using its proved
regularity; it does not presume that $f$ is injective.
\end{proof}
% END WIS12

% BEGIN WIS13
\paragraph{WIS13: injectivity and all $p$-power saturations.}
\label{claim:WIS13}
The map $f$ is injective. For all $r\geq0$ and $x\in E$,
\[
 f(x)\in p^rW\quad\Longrightarrow\quad x\in p^rE.
                                                               \tag{WIS.19}
\]
If $f(x)=p^rw$, there is $x'\in E$ with
$x=p^rx'$ and $f(x')=w$.

\begin{proof}
If $f(x)=0$, WIS12 makes $x$ divisible by every
$t_E^{er}$. Since $p^r=u_E^rt_E^{er}$ with $u_E$
a unit, $x$ lies in every $p^rE$. Separatedness gives
$x=0$, proving injectivity.

For (WIS.19), the equality $p^r=u_W^rt_W^{er}$
shows that $t_W^{er}$ divides $f(x)$. Apply WIS12
and the same equality in $E$. Choose $x'$ with
$x=p^rx'$. The relation $p^r(f(x')-w)=0$ and
the absence of $p$-torsion in $W$ give $f(x')=w$.
This includes $r=0$.
\end{proof}
% END WIS13

% BEGIN WIS14
\paragraph{WIS14: reflection of units.}
\label{claim:WIS14}
If $f(x)$ is a unit of $W$, then $x$ is a unit of $E$.

\begin{proof}
Suppose $x$ is not a unit, and choose a maximal ideal
$M$ of $E$ containing $x$. By WIS10, $t_E\in M$.
Choose (WIS.16); it gives $j(a)\in M$.

For any $i$, the first alternative $e_i a=tb$ in
(WIS.11) is impossible. Indeed, multiplying (WIS.16)
by $j(e_i)$ and mapping to $W$ would give
\[
 f(j(e_i)) f(x)\pi_W^m
       =t_W\bigl(f(j(b))+f(j(e_i))f(z)\bigr).
 \tag{WIS.20}
\]
The factor $f(x)\pi_W^m$ is a unit, contradicting WIS10.
Hence the second alternative holds for every $i$.
Its image in $E$ is
\[
         \pi_E^kv_Ej(e_i)=j(e_i)j(a)-t_Ej(b)\in M.
\]
Canceling the unit on the left gives $j(e_i)\in M$
for every $i$. Their sum is one, contradicting that
$M$ is proper. Therefore $x$ is a unit.
\end{proof}
% END WIS14

\subsection{The integral intersection}

% BEGIN WIS15
\paragraph{WIS15: removing one integral uniformizer digit.}
\label{claim:WIS15}
For $x\in E$ with $f(x)\in c(B)$, there are $a\in A$
and $x'\in E$ such that
\[
                   x=j(a)+t_Ex',\qquad f(x')\in c(B).      \tag{WIS.21}
\]

\begin{proof}
Write $f(x)=c(y)$ and choose (WIS.16).
We first prove, for every $i$, an expression
\[
                         e_i a=tb_i+\pi^md_i.             \tag{WIS.22}
\]
The first alternative in (WIS.11) gives this with $d_i=0$.
In the second alternative write
$e_i a=tb_i+\pi^{k_i}v_i e_i$ with $v_i$ a unit.
If $k_i\geq m$, take $d_i=\pi^{k_i-m}v_i e_i$.

Suppose instead that $k_i<m$, and put $h_i=f(j(e_i))$.
Multiplying (WIS.16) by the component idempotent and
using $f(x)=c(y)$ gives
\[
 \pi_W^{k_i}(v_i)_W h_i
    \left(1-c(y)\pi_W^{m-k_i}(v_i)_W^{-1}\right)
        =-t_W\bigl(f(j(b_i))+h_i f(z)\bigr).                \tag{WIS.23}
\]
The factor in parentheses is
\[
          c\left(1-y\,g(\pi)^{m-k_i}g(v_i)^{-1}\right).
\]
Because $m-k_i\geq1$ and $g(\pi)\in\operatorname{Jac}(B)$,
this factor is a unit. The two other factors preceding
$h_i$ in (WIS.23) are units as well. Consequently $t_W$
would divide $h_i$, contradicting WIS10. This proves
(WIS.22) in every case.

Summing (WIS.22) gives $a=tb+\pi^md$.
Substitution into (WIS.16), followed by cancellation
of $\pi_E^m$, gives
\[
         x=j(d)+t_Ex',\qquad
         x'=(j(b)+z)\pi_E^{-m}.                            \tag{WIS.24}
\]
Its image says $c(y-g(d))=t_W f(x')$. Multiply by
$u_Wt_W^{e-1}$ and use $p=u_Wt_W^e$:
\[
       c\bigl(g(u)g(t)^{e-1}(y-g(d))\bigr)=p f(x').
\]
The $p$-saturation of $c(B)$ now proves $f(x')\in c(B)$.
This is the second assertion in (WIS.21), which is
essential for iterating the construction.
\end{proof}
% END WIS15

% BEGIN WIS16
\paragraph{WIS16: the exact integral intersection.}
\label{claim:WIS16}
Under the abstract hypotheses,
\[
             f(x)\in c(B)\quad\Longrightarrow\quad
                          x=j(a)\text{ for some }a\in A. \tag{WIS.25}
\]
Consequently
\[
                       f(E)\cap c(B)=f(j(A))=c(g(A)).     \tag{WIS.26}
\]

\begin{proof}
Iterate WIS15 finitely many times. At every integer $N$
it gives an $a^{(N)}\in A$ such that
$x-j(a^{(N)})\in t_E^NE$. Take $N=ek$ and use
$p^k=u_E^kt_E^{ek}$ to choose $a_k$ with
\[
                         x-j(a_k)\in p^kE.               \tag{WIS.27}
\]
For $k\leq k'$, subtraction shows
$j(a_k-a_{k'})\in p^kE$. The stated saturation of
$j:A\to E$ implies $a_k-a_{k'}\in p^kA$.
Completeness (WIS.12) supplies $a\in A$ with
$a_k-a\in p^kA$ for all $k$.

Maps of rings preserve divisibility by the integer $p^k$,
so $j(a_k)-j(a)\in p^kE$. Adding this to (WIS.27)
puts $x-j(a)$ in every $p^kE$. Separatedness of $E$
gives $x=j(a)$.
No continuity of $j$ for an additional topology is
required: the relevant congruences are literal ring
divisibilities. The reverse containment in (WIS.26)
is the commutative-square identity $fj=cg$.
\end{proof}
% END WIS16

% BEGIN WIS17
\paragraph{WIS17: the full abstract four-clause embedding theorem.}
\label{claim:WIS17}
For the square (WIS.14) with the hypotheses printed after it,
$f$ is injective, reflects units, satisfies
\[
 f(x)=py\quad\Longrightarrow\quad x=px'
                                  \text{ for some }x'\in E,
\]
and satisfies the intersection assertion (WIS.25).
These are exactly the four conclusions of the original
abstract theorem in the Wach-inclusion proof.

\begin{proof}
Injectivity and the displayed saturation are WIS13,
unit reflection is WIS14, and the last assertion is WIS16.
Their proofs use the listed finite family of idempotents,
its dichotomy and nonvanishing, not any additional
orthogonality, basis, or domain hypothesis on $E,B,W$.
The completeness assumption on $A$ is precisely the
existence of limits for (WIS.12); the proof did not
assume its separatedness.
\end{proof}
% END WIS17

% BEGIN WIS18
\paragraph{WIS18: the full original Wach-inclusion theorem.}
\label{claim:WIS18}
Fix every original allowed tuple
$(p,P,AI,K,s,L)$, its condition $\mathrm{Good}(K,s,L)$,
and the supplied records $\Gamma,R,I$ with $\mathrm{IoW}(I)$.
Set $O=\mathcal O_L$, $n=[K_0:\mathbb Q_p]$,
$A=\prod_{\mathbb Z/n\mathbb Z}O[[T]]$, $E=R.EE$,
and let $B,W,c$ be the actual rings and map (WIS.7).
With the original maps
\[
 j:A\longrightarrow E,\qquad g=I.io:A\longrightarrow B,
                   \qquad f=I.ioE:E\longrightarrow W,
\]
all four original conclusions hold:
\begin{enumerate}
\item $f$ is injective.
\item If $f(x)$ is a unit, then $x$ is a unit.
\item If $f(x)=py$ in $W$, then $x=px'$ for some $x'\in E$.
\item If $f(x)=c(y)$ for some $y\in B$, then $x=j(a)$
      for some $a\in A$.
\end{enumerate}

\begin{proof}
We check every input of WIS17 on the actual data.
By WIS1 choose $\varpi,u,e$ with $p=u\varpi^e$.
In $A$ take the constant tuples $t=(\varpi)_i$ and
$u_A=(u)_i$, and $\pi=(T)_i$.
The original period datum identifies its residue
characteristic with $p$; thus the completeness,
density, saturation and invertibility clauses of
$R$ are precisely
\[
 \begin{gathered}
 E\text{ is }p\text{-adically complete and separated},
 \qquad j(\pi)\in E^\times,\\
 \forall x\in E,\ \exists a,m:\ xj(\pi)^m-j(a)\in pE,\\
 j(a)\in p^kE\Longrightarrow a\in p^kA .
 \end{gathered}                                            \tag{WIS.28}
\]
The density clause here is its original clause at
precision one; no stronger description of $E$ is used.
The original $I.ioE_{\mathrm{ext}}$ is $fj=cg$, and
$I.io_\pi$ identifies $g(\pi)$ with $1\otimes\mu_s$.

The component idempotents sum to one and WIS7 gives
their dichotomy. WIS8 gives completeness of $A$.
The integer $n$ is positive because $K_0$ is a finite
field extension of $\mathbb Q_p$. Its component
Frobenius shifts the indices, and the two original
Frobenius-compatibility clauses give (WIS.13).
WIS5 makes $W$ nonzero, so WIS9 gives nonvanishing
of every component image.

Finally WIS5 supplies absence of $p$-torsion and
$p$-saturation of $c(B)$, and WIS6 supplies both
required Jacobson-radical assertions. WIS17 therefore
applies to the exact original maps, proving all four
clauses. The hypotheses $\mathrm{Good}$ and
$\mathrm{IoW}$ are retained at the original endpoint;
this particular proof does not need their additional
arithmetic consequences once the displayed records
have been supplied.
\end{proof}
% END WIS18

\subsection{Closed images and fixed elements}

% BEGIN WIS19
\paragraph{WIS19: the original $p$-adic closed-image theorem.}
\label{claim:WIS19}
Let $j_0:E_0\to R_0$ be a unital map of commutative rings.
Assume $R_0$ is $p$-torsion-free and $p$-adically separated,
$E_0$ is $p$-adically complete and separated, and
\[
 j_0(e)=pr\Longrightarrow e=pe'
                              \text{ for some }e'\in E_0. \tag{WIS.29}
\]
If $x\in R_0$ admits approximations
$x-j_0(e_k)\in p^kR_0$ for every $k$, then
$x=j_0(e)$ for some $e\in E_0$.
In particular the actual image $I.ioE(E)$ is $p$-adically
closed in $W$, and the induced map
$E/p^kE\to W/p^kW$ is injective for every $k$.

\begin{proof}
First iterate (WIS.29). If $j_0(e)=p^{k+1}r$,
write $e=pe'$ using (WIS.29). Cancellation of $p$
in $R_0$ gives $j_0(e')=p^kr$, and induction gives
$e\in p^{k+1}E_0$. Thus $j_0$ reflects every
$p$-power divisibility, without an assumption that
$j_0$ is injective.

For $k\leq k'$ the approximations imply
$j_0(e_k-e_{k'})\in p^kR_0$. The reflected
divisibility makes $(e_k)$ a $p$-adic Cauchy system.
Choose its limit $e$. Both
$x-j_0(e_k)$ and $j_0(e_k-e)$ lie in $p^kR_0$,
so $x-j_0(e)$ lies in every $p^kR_0$ and is zero.
This proves the full original generic
\texttt{padic\_closed} implication.

Apply it to $I.ioE$ using (WIS.28), WIS5 and WIS18.
The kernel of the map modulo $p^k$ is zero by
(WIS.19), proving the last assertion as well.
\end{proof}
% END WIS19

% BEGIN WIS20
\paragraph{WIS20: the full fixed-subring $p$-adic supplier.}
\label{claim:WIS20}
For any set $H$ of elements of the original Galois group, let
\[
             W^H=\{w\in W:\sigma(w)=w\text{ for all }\sigma\in H\}.
\]
No subgroup or closedness hypothesis on $H$ is required.
The following original properties hold:
\begin{enumerate}
\item $pw\in W^H$ implies $w\in W^H$.
\item $w\in W^H$ implies $\phi(w)\in W^H$.
\item If $w_i\in W^H$ for every $i$, there exists $w\in W^H$
      such that for every $k$,
      \[
                  w-\sum_{i<k}p^iw_i\in p^kW.             \tag{WIS.30}
      \]
\end{enumerate}

\begin{proof}
The Galois action is the identity on $O$ and the
Witt-vector lift of its action on $F$.
It fixes $p$. Hence $pw$ fixed gives
$p(\sigma(w)-w)=0$, and WIS5 permits cancellation.

On $F$ every ring automorphism commutes with $x\mapsto x^p$.
Applying this equality in Witt coordinates and then
tensoring with the identity on $O$ shows
$\sigma\phi=\phi\sigma$. This proves the second clause.

The partial sums in (WIS.30) form a $p$-adic Cauchy
system in the complete ring $W$, so they have a
limit $w$ satisfying its stated congruences.
For every $\sigma\in H$, applying $\sigma$ to each
congruence is legitimate algebraically: a ring map
preserves $p^kW$. Since the partial sum is fixed,
$\sigma(w)-w\in p^kW$ for every $k$.
Separatedness of $W$ gives $\sigma(w)=w$.
This proves the full original
\texttt{fixSub\_ladic} statement, including its
summation conclusion for arbitrary $H$.
\end{proof}
% END WIS20

% BEGIN WIS21
\paragraph{WIS21: an actual fixed lift of an element fixed modulo $p$.}
\label{claim:WIS21}
For any set $H$ as above, suppose $x\in W$ satisfies
$\sigma(x)-x\in pW$ for all $\sigma\in H$.
Then there is $y\in W^H$ with $y-x\in pW$.
Using the same integral basis as WIS4, it can be chosen as
\[
            y=\sum_i b_i\otimes[(x_i)_0],
            \qquad x=\sum_i b_i\otimes x_i.                \tag{WIS.31}
\]

\begin{proof}
By WIS4 and WIS2, the congruence for $\sigma(x)-x$
is equivalent to
$\sigma((x_i)_0)=(x_i)_0$ in $F$ for each $i$.
The Teichmuller lift commutes with every coefficient
automorphism: $\sigma([a])=[\sigma(a)]$.
Thus each term in (WIS.31) is fixed and $y\in W^H$.
Its coordinates have the same zeroth Witt coordinates
as those of $x$, so $y-x\in pW$ by WIS2 and WIS4.
The finite coefficient basis is fixed by the original
action; averaging over $H$ is neither needed nor used.
This proves the full original fixed-modulo-prime
lifting statement for every permitted set $H$.
\end{proof}
% END WIS21

\subsection{The exact cyclotomic approximation filtration}

For the rest of this subsection keep the original $P,AI,L,s$,
but no series record $R$ or embedding $I$ is required.
Put
\[
 \tau=\epsilon_s-1\in S,\qquad
 \pi=c(1\otimes\mu_s)\in W,\qquad
 \mathcal N_m=\pi^mc(B)+pW\quad(m\geq0).                    \tag{WIS.32}
\]
Each $\mathcal N_m$ is an additive subgroup. It is not being
treated as the ideal $\pi^mW+pW$; in the embedding situation
$\pi$ is a unit of $W$, so that ideal would lose the
filtration. The restriction of the coefficient to $c(B)$
in (WIS.32) is essential.

% BEGIN WIS22
\paragraph{WIS22: separation and complete summation in the actual tilt.}
\label{claim:WIS22}
For every $s\geq0$, the original tilt ring $S$ satisfies
\[
 \bigcap_{m\geq0}\tau^mS=0,\qquad
 \forall(a_i)_{i\geq0}\subset S,\quad
 \exists a\in S:\quad
 a-\sum_{i<m}\tau^ia_i\in\tau^mS\quad\text{for every }m.
                                                               \tag{WIS.33}
\]
These are the two clauses of the original
tilt separation-and-completeness theorem.

\begin{proof}
In the inverse-limit description of $S$, the $r$th
coordinate of $\epsilon_s$ is
$\zeta_{s+r}\bmod p$: inverse Frobenius shifts the
compatible coordinates. The cyclotomic identities give
$\zeta_{s+r}^{p^{s+r}}=1$. Thus, in $\mathcal O_C/p$,
\[
                   \tau_r^{M(r)}=0,\qquad M(r)=p^{s+r}.
                                                               \tag{WIS.34}
\]
If $a$ is divisible by every power of $\tau$, use
$\tau^{M(r)}\mid a$ to conclude that its $r$th
coordinate is zero. This for all $r$ proves separation.

For a sequence $(a_i)$, the $r$th coordinate of
$\sum_{i<N}\tau^ia_i$ is constant once $N\geq M(r)$.
Define $a_r$ to be that stable value. These coordinates
are compatible under Frobenius: compare adjacent
coordinates at any common cutoff at least
$\max(M(r),M(r+1))$, and use compatibility of the
finite sum. Hence they define $a\in S$.

For a fixed $m$, repeat this construction for the tail
sequence $a_m,a_{m+1},\ldots$, obtaining $b_m\in S$.
At each coordinate, evaluate the two sums at a common
large finite cutoff. Cancellation of the first $m$
terms then gives
\[
                     a-\sum_{i<m}\tau^ia_i=\tau^m b_m.
 \tag{WIS.35}
\]
This proves the complete summation assertion, including
literal divisibility of every tail. No convergence of
a chosen sequence in the characteristic-zero period
ring is assumed here; all limits were constructed
by stabilized finite coordinates in the tilt.
\end{proof}
% END WIS22

% BEGIN WIS23
\paragraph{WIS23: the original approximation condition in coordinates.}
\label{claim:WIS23}
For $x=\sum_i b_i\otimes x_i\in W$, write
$q_i(x)=(x_i)_0\in F$. Let $\iota:S\hookrightarrow F$
be the original inclusion. Then
\[
 \begin{split}
 x\in pW&\quad\Longleftrightarrow\quad q_i(x)=0
                                      \text{ for all }i,\\
 x\in\mathcal N_m&\quad\Longleftrightarrow\quad
 q_i(x)\in\iota(\tau^mS)\text{ for all }i.
 \end{split}                                               \tag{WIS.36}
\]
These are the actual \texttt{dvd\_iff\_crd} and
\texttt{nearPi\_iff\_crd} criteria, with their prescribed
maps and parameter.

\begin{proof}
The first assertion follows from WIS2 and WIS4.
The zeroth Witt coordinate of $\mu_s$ is $\tau$.
Hence the coordinates of $\pi^mc(b)$ at degree zero
are $\iota(\tau^m d_i)$, where $d_i$ are the zeroth
coordinates of $b$. The term $pw$ contributes zero,
proving the forward implication of the second assertion.

Conversely, write $q_i(x)=\iota(\tau^m d_i)$.
Choose $b=\sum_i b_i\otimes[d_i]\in B$.
The degree-zero coordinates of $x-\pi^mc(b)$
are all zero, so this difference equals $pw$
for some $w\in W$. This is precisely the original
definition of membership in $\mathcal N_m$.
All coordinate maps and inclusions used here are
the actual tensor maps, not independently selected
identifications of quotient rings.
\end{proof}
% END WIS23

% BEGIN WIS24
\paragraph{WIS24: the full original approximation separation and summation.}
\label{claim:WIS24}
For the subgroups (WIS.32),
\[
       \bigcap_{m\geq0}\mathcal N_m=pW.                    \tag{WIS.37}
\]
If $k\geq0$ and $f_i\in\mathcal N_{i+k}$ for every $i$,
there is $x\in W$ such that
\[
                      x-\sum_{i<m}f_i\in\mathcal N_{m+k}
                          \quad\text{for all }m\geq0.     \tag{WIS.38}
\]
This proves both clauses of the original
\texttt{nearPi\_sep\_complete} theorem for every $s$.

\begin{proof}
Certainly $pW$ lies in every $\mathcal N_m$.
If $x$ lies in their intersection, (WIS.36) at $m=0$
writes $q_i(x)=\iota(d_i)$ with $d_i\in S$.
At any other $m$, the injectivity of $\iota$ says
$d_i\in\tau^mS$. WIS22 makes each $d_i$ zero,
and (WIS.36) gives $x\in pW$.

For (WIS.38), choose $d_{i,j}\in S$ with
\[
                         q_j(f_i)=\iota(\tau^{i+k}d_{i,j}).
\]
For each of the finitely many coefficient indices $j$,
WIS22 constructs $a_j\in S$ such that
$a_j-\sum_{i<m}\tau^id_{i,j}\in\tau^mS$.
Take the actual element
\[
                      x=\sum_j b_j\otimes[\iota(\tau^k a_j)]
                               \in W.
 \tag{WIS.39}
\]
Its difference from the partial sum has every
degree-zero coordinate in $\iota(\tau^{m+k}S)$.
WIS23 gives (WIS.38). This proves the exact
approximation statement; it does not replace it
with an assertion about a different topology on $W$.
\end{proof}
% END WIS24

% BEGIN WIS25
\paragraph{WIS25: the full original fixed-subring approximation supplier.}
\label{claim:WIS25}
Let $H$ be any set of elements of the original Galois group
and suppose $\pi\in W^H$. If $f_i\in W^H\cap
\mathcal N_{i+k}$ for every $i$, there exists $y\in W^H$
such that
\[
                    y-\sum_{i<m}f_i\in\mathcal N_{m+k}
                               \quad\text{for every }m.  \tag{WIS.40}
\]
Thus the complete original
\texttt{fixSub\_nearPi\_complete} conclusion holds.

\begin{proof}
Apply WIS24 first without a fixedness requirement
and obtain $x$ satisfying (WIS.38). Each original
Galois action preserves $c(B)$, because it acts
on the tilt and its Witt vectors and fixes $O$.
For $\sigma\in H$, it also fixes $\pi$. Thus it
preserves each $\mathcal N_m$.

Let $S_m=\sum_{i<m}f_i$. This partial sum is fixed, so
\[
              \sigma(x)-x
                =\sigma(x-S_m)-(x-S_m)\in\mathcal N_{m+k}.
 \tag{WIS.41}
\]
Since $\pi\in c(B)$, the subgroups $\mathcal N_m$
decrease with $m$. It follows that the left side
of (WIS.41) lies in every $\mathcal N_m$.
By WIS24 it lies in $pW$. WIS21 therefore supplies
an actual fixed $y$ with $y-x\in pW$.
As $pW\subseteq\mathcal N_{m+k}$ for every $m$,
adding this correction to (WIS.38) proves (WIS.40).
The same basis and the same parameter $\pi$ have
been used throughout. In particular no additional
restriction on $H$ or on the allowed index $s$
has been imposed.
\end{proof}
% END WIS25

% BEGIN WIS26
\paragraph{WIS26: the precise return to the Wach-lattice construction.}
\label{claim:WIS26}
On every original permitted layer with its supplied
$\Gamma,R,I,E_{\mathrm{et}}$, retain the original
$\operatorname{LocTied}(p,P,AI)$ hypothesis (hence the
$\operatorname{IsStd}$ instance used for realization), together
with the original Bhatt--Scholze crystalline-realization input. Suppose
the following two original suppliers are available:
\begin{enumerate}
\item The actual Wach prism with its maps, coefficient
      coordinates and scalar-map identities.
\item Descent of the original cyclotomic-kernel-fixed
      homogeneous Frobenius matrix solutions, at every
      required rank on that layer.
\end{enumerate}
Then the full original Wach-lattice existence conclusion
of WLD26 holds. A separate assumption of injectivity or
unit reflection for $I.ioE$ is no longer needed.

\begin{proof}
Apply WIS18 to the supplied $I$ to obtain the two
clauses used by WLD26, namely injectivity and unit
reflection. The two retained suppliers are exactly
the other two period-ring inputs in that theorem.
All its original coefficient data and realizing
crystals remain unchanged. Thus WLD26 applies,
producing the same inverse-image lattice
$N_T=j_T^{-1}(L_T)$ inside the supplied
\'etale module, with all the fields already
proved in WLD18--WLD25.

This implication is uniform over the original
allowed layer indices and over the supplied
\'etale records. No transition maps between
different outer indices are added. The independent
faithful-invariant, \'etale-basis, Herr-class-map
and quantitative comparison inputs needed after
Wach existence are not consequences of this
two-supplier return.
\end{proof}
% END WIS26

\subsection{The remaining period-ring work}

The embedding theorem WIS18 has its full four-clause original
scope. Its first two clauses feed the lattice construction;
its saturation clause also feeds the original invariant
Frobenius-matrix proof. WIS19 supplies that proof's
$p$-adic closed-image input, WIS20 its fixed-subring
$p$-adic input, and WIS23--WIS25 its original approximation
separation and completeness inputs.

The matrix-solution theorem itself is still to be reconstructed.
Its remaining argument involves the Frobenius behavior of the
approximation filtration, approximation of actual cyclotomic
invariants by the specified coefficients, the closedness
modulo $p$ in that filtration, and the matrix lifting step.
The exact source endpoint is
\texttt{FixedByCyclotomicKernelEtalePhiSolutionMemWach}.
These requirements are not supplied merely by the existence
of the limits in WIS24 or WIS25.

The actual prism in
\texttt{ExistsWachPrismOverIntegers} is the other remaining
period-ring supplier. Both have concrete maps and identities
that must be verified before WIS26 becomes an unconditional
Wach-existence result. The later full local comparison package
and the common family of quantitative local bounds are
separate subsequent tasks. No whole selected global obligation
is declared closed in this chapter.


